Die Erhöhung der Resilienz dezentraler Ladeinfrastrukturen baut auf bestehenden Erkenntnissen der verteilten Angriffserkennung und Informationsfusion auf. Die folgenden Abschnitte fassen den aktuellen Stand der Forschung in den Bereichen der Alert-Aggregation, der permutationsinvarianten Mengenverarbeitung sowie kollaborativer Sharing-Strategien zusammen.

\section{Anomalieerkennung und Alert-Aggregation in verteilten Systemen}
Lokale Detektionsverfahren stoßen strukturell an Grenzen, wenn netzwerkweite oder koordinierte Angriffe vorliegen, die auf Ebene eines einzelnen Knotens nicht als bösartig identifiziert werden können. Lo et al. adressieren dieses Problem mit E-GraphSAGE \cite{lo2022egraphsage}, einem GNN-basierten IDS für IoT-Netzwerke, das topologische Informationen und Kanten-Features aus Netzwerkflussdaten für die Angriffserkennung kombiniert. Heigl et al. präsentieren mit dem SOAAPR-Framework \cite{heigl2021exploiting} einen Ansatz, der Ausreißererkennungsergebnisse auf Streaming-Daten auswertet und daraus neuartige Angriffsmuster extrahiert. Boswell et al. schlagen mit FLARE \cite{boswell2025flare} eine leichtgewichtige Feature-Aggregationstechnik für IoT-IDS vor, die mittels zeitbasierter Sliding Windows statistische Kennzahlen wie Flussrate, Paketgrößen und Entropie aggregiert, um kurzlebige Angriffsmuster wie DoS und DDoS zuverlässig zu erkennen.

Auf Ebene der Alert-Konsolidierung beschreiben Debar \& Wespi \cite{debar2001aggregation} einen grundlegenden\\Aggregations- und Korrelationsalgorithmus, der Duplikate und Folgewarnungen erkennt sowie Alerts situationsbezogen zusammenfasst. Zhang et al. \cite{zhang2019iacf} führen diesen Gedanken mit dem IACF-Framework fort, indem sie Alarme basierend auf ihren inneren Zusammenhängen als \textit{Intrusion Actions} modellieren. Durch den Einsatz einer zeitbasierten Sequenzanalyse lassen sich so komplexe, mehrstufige Angriffsmuster sowohl rekonstruieren als auch deren künftigen Verlauf vorhersagen. Im Gegensatz zu diesen regelbasierten Verfahren verfolgt die vorliegende Arbeit einen lernbasierten, gewichteten Aggregationsansatz, der die Relevanz einzelner Threat-Reports dynamisch auf Basis enthaltener Einzeldetektionen bewertet.

\section{Informationsfusion mit variabler Eingabegröße und Attention-Mechanismen}
Eine zentrale Herausforderung bei der Aggregation verteilter Threat-Reports ist die variable Anzahl eingehender Quellen, da klassische neuronale Netze (Fixed-Size Input) erfordern. Zaheer et al. lösen dieses Problem mit der formalen Charakterisierung permutationsinvarianter Funktionen auf Mengen (Deep Sets) \cite{zaheer2017deepsets}: Jede solche Funktion lässt sich als Summen-Dekomposition darstellen, wobei Sölch et al. empirisch zeigen \cite{soelch2019deepsetAgg}, dass die Wahl der Aggregationsfunktion die Modellperformance erheblich beeinflusst und lernbare Mechanismen statischen Operatoren wie Mean oder Max überlegen sein können. Vinyals et al. demonstrieren mit Set2Set \cite{vinyals2016orderm} durch algorithmische Experimente, dass eine reihenfolgeunabhängige, iterative Verarbeitung von Mengen die Lern- und Generalisierungsfähigkeit von Modellen signifikant steigert, Lee et al. bauen darauf mit dem Set Transformer \cite{lee2019settransformer} auf, der Self-Attention zur Kodierung von Element-Interaktionen nutzt und dabei die quadratische Komplexität klassischer Attention auf lineare Komplexität reduziert.

Als Grundlage für die in dieser Arbeit verwendeten Attention- und Masking-Mechanismen dienen Vaswani et al. \cite{vaswani2017attention}, die die Transformer-Architektur einführen, in der Padding und Masking die effiziente Batch-Verarbeitung variabel langer Sequenzen ermöglichen. Veličković et al. übertragen dieses Prinzip mit Graph Attention Networks (GATs) \cite{velickovic2018gat} auf graphbasierte Strukturen und nutzen Masked Self-Attention zur Berechnung von Knoten-Repräsentationen ohne kostspielige Matrixoperationen. Liu et al. \cite{liu2022pooling} liefern einen systematischen Überblick über Pooling-Operationen im Deep Learning und bilden damit die kategorische Einordnung der in dieser Arbeit evaluierten Aggregationsmethoden.

\section{Kollaboratives Threat-Report-Sharing, Gewichtungsstrategien und Datenqualität}
Das kollaborative Teilen von Threat-Reports zwischen Knoten erfordert Mechanismen zur Bewertung der Vertrauenswürdigkeit einzelner Knoten. Zhang, Miao \& Xue \cite{zhang2022reputation} schlagen ein blockchain-basiertes Modell vor, das ein verteiltes Reputationsmanagementsystem integriert, dabei zeigen sie, dass ein Ungleichgewicht zwischen normalen und kompromittierten Knoten die Aggregation systematisch verzerren kann, ein Phänomen, das in der vorliegenden Arbeit als Echo-Chamber-Effekt formalisiert und durch asymptotisches Sättigungsverhalten der Gewichtungsfunktionen strukturell begrenzt wird. Zur Modellierung dieses abnehmenden Grenznutzens greift die Arbeit auf die von Norstad \cite{norstad1999utility} beschriebene Negative Exponential Utility Function mit konstanter absoluter Risikoaversion (CARA) zurück, sowie auf die von Pérez-Maldonado et al. \cite{perezmaldonado2025sigmoid} formal hergeleitete logistische Sigmoid-Funktion als S-förmige Alternative.

Hinsichtlich Datenqualität zeigen Hasan et al. \cite{hasan2021missingvalue} in einem systematischen Review, dass fehlende Werte die Performance von ML-Modellen erheblich beeinflussen, eine Grundlage für die implementierte NaN-Fehlerbehandlung in der Threat-Report-Generierung. Die Entscheidung für das CSV- gegenüber dem Parquet-Format wird durch Analysen von Taylor-Davies \cite{taylordavies2023parquet} sowie CelerData \cite{celerdata2024parquet} gestützt, die belegen, dass der Metadaten-Overhead von Parquet bei kleinen Dateien unter 30~kB die Vorteile der spaltenorientierten Kompression überwiegt.